\section{Probabilidade Condicional e Independência}

A probabilidade condicional representa o conceito de probabilidade de um evento tal que
é conhecida a ocorrência de outro; como se tivéssemos reduzido o espaço amostral a um
conjunto menor de resultados possíveis.

\begin{def:probabilidade condicional}
Sejam $E$ e $F$ eventos num espaço amostral $\Omega$. A probabilidade condicional
do evento $E$ dado $F$, denotada por $P(E|F)$, é dada por
\[
	P(E|F) = \frac{P(E \cap F)}{P(F)}
\]
\end{def:probabilidade condicional}

A probablidade condicional induz uma forma de determinar a probabilidade do evento
dado como interseção de vários outros, expressão essa denominada \textbf{regra do produto}.
\begin{prop:regra da multiplicacao}
Se $\mathcal{F} = \{E_1, \ldots, E_n\}$ uma família de eventos de $\Omega$, então
\[
	P\left(\bigcap_{E \in \mathcal{F}} E\right) = \prod_{i=1}^n P\left(E_i\ \middle|\ \bigcap_{j < i} E\right)
\]
\begin{proof}
	A prova será por indução em $n$

	Base: $n = 1$ é trivial.

	Hipótse de Indução: admita que o resultado vale para uma família de $n$ eventos. Se
	$|\mathcal{F}| = n + 1$, então façamos
	\[
		\bigcap_{j < n+1} E_j  = E'
	\]
	e obtemos que
	\[
		P\left(\bigcap_{E \in \mathcal{F}} E\right) = P(E' \cap E_n) = P(E_n | E') P(E').
	\]
	Como a hipótese de indução vale para $E'$, então
	\[
		P\left(\bigcap_{E \in \mathcal{F}} E\right) = P\left(E_n\ \middle|\ \bigcap_{j < n+1} E_j\right)
		\prod_{i=1}^n P\left(E_i\ \middle|\ \bigcap_{j < i} E\right) =
		\prod_{i=1}^{n+1} P\left(E_i\ \middle|\ \bigcap_{j < i} E\right)
	\]
	Uma vez que a proposição vale para $n+1$, então vale para todo $n \in \mathbb{N}$
	pelo princípio de indução.
\end{proof}
\end{prop:regra da multiplicacao}

Eventos independentes são aqueles cuja probabilidade de ocorrência não depende
da ocorrência de outro. Isso equivale a dizer que, o conhecimento da ocorrência
de evento não muda a crença da ocorrência de outro.
\begin{def:independencia}
Dois eventos $E$ e $F$ dum espaço amostral $\Omega$ são independentes se, e somente
se
\[
	P(A | B) = P(A) \cdot P(B)
\]
donde se conclui que
\[
	P(A | B) = P(A) \quad \mathit{ e }\quad P(B | A) = P(B)
\]
\end{def:independencia}

\begin{prop:independencia}
Sejam $E$ e $F$ eventos em $\Omega$. Se eles são independentes, então os seguintes
eventos também o são:
\begin{enumerate}
	\item $E^c$ e $F$
	\item $E$ e $F^c$
	\item $E^c$ e $F^c$
\end{enumerate}
\begin{proof}
	(1) É verdade que
	\[
		(E \cap F) \cap (E \cap F^c) = \varnothing
	\]
	e também que
	\[
		(E \cap F) \cup (E \cap F^c) =  E
	\]
	então
	\[
		P(E \cap F) + P(E \cap F^c) = P(E)
	\]
	donde segue que
	\begin{gather*}
		P(E \cap F^c) = P(E) - P(E \cap F) = P(E) - P(E)P(F) = P(E) \cdot [1 - P(F)] \\
		P(E \cap F^c) = P(E) \cdot P(F^c)
	\end{gather*}
	A expressão (2) é análoga a (1). Por fim, para expressão (3), conclua que
	\begin{align*}
		P(E^c \cap F^c) & = P((E \cup F)^c) = 1 - P(E \cup F) \\
		                & = 1 - P(E) - P(F) + P(E \cap F)     \\
		                & = P(E^c) - P(F) + P(E) \cdot P(F)   \\
		                & = P(E^c) - P(F) [1 - P(E)]          \\
		                & = P(E^c) - P(F) \cdot P(E^c)        \\
		                & = P(E^c) [1 - P(F)]                 \\
		                & = P(E^c) \cdot P(F^c)               \\
	\end{align*}
\end{proof}
\end{prop:independencia}
